\documentclass[11pt]{article}
\usepackage{amsmath,amsthm,amssymb}
\usepackage[margin=1in]{geometry}

\newtheorem{theorem}{Theorem}
\newtheorem{lemma}[theorem]{Lemma}
\newtheorem{proposition}[theorem]{Proposition}
\newtheorem{corollary}[theorem]{Corollary}
\theoremstyle{definition}
\newtheorem{definition}[theorem]{Definition}
\theoremstyle{remark}
\newtheorem{remark}[theorem]{Remark}

\DeclareMathOperator{\rank}{rank}
\DeclareMathOperator{\tr}{tr}
\DeclareMathOperator{\im}{im}
\DeclareMathOperator{\sgn}{sgn}

\title{The H-Invariant Theorem for Staircase Products\\in the Symmetric Group}
\author{Clio}
\date{April 17, 2026}

\begin{document}
\maketitle

\begin{abstract}
Let $w_0 \in S_n$ be the longest element and let
$\Pi = \prod_{j=1}^{m}(1+s_{i_j})$ be the staircase product, where
$s_{i_1},\ldots,s_{i_m}$ is the staircase reduced word
$s_1,\,s_2 s_1,\,s_3 s_2 s_1,\ldots,s_{n-1}\cdots s_1$ for~$w_0$.
Let $H = \langle s_1, s_3, s_5, \ldots\rangle$ be the Young subgroup
$S_{(2,\ldots,2,1,\ldots)}$. We prove that
\[
  \rank(\Pi|_{V_\lambda}) \;=\; \dim(V_\lambda^H)
\]
for every irreducible representation~$V_\lambda$ of~$S_n$.
The proof combines W-graph positivity (giving non-negative traces),
exterior powers (reducing the lower bound to non-vanishing of traces on
specific irreducibles), and the seminormal representation (establishing
the required non-vanishing via column-adjacency-free standard Young tableaux).
\end{abstract}

%% =====================================================================
\section{Setup and notation}
%% =====================================================================

Fix $n \ge 2$. The symmetric group $S_n$ is generated by simple
transpositions $s_1,\ldots,s_{n-1}$, where $s_i = (i\;\;i{+}1)$.
The \emph{staircase reduced word} for the longest element
$w_0 \in S_n$ is
\[
  \mathbf{i} \;=\; (s_1),\;(s_2,s_1),\;(s_3,s_2,s_1),\;\ldots,\;
  (s_{n-1},\ldots,s_2,s_1),
\]
a sequence of $m = \binom{n}{2}$ simple transpositions.
Define the \emph{staircase product}
\[
  \Pi \;=\; \prod_{j=1}^{m}(1+s_{i_j})
  \;\in\; \mathbb{Q}[S_n].
\]
Let $H = \langle s_1, s_3, s_5, \ldots \rangle \cong
S_2 \times S_2 \times \cdots$, the Young subgroup corresponding to the
composition $(2,2,\ldots,2,1,\ldots,1)$ of~$n$.
For an irreducible representation $V_\lambda$ indexed by
$\lambda \vdash n$, write
$r_\lambda = \rank(\Pi|_{V_\lambda})$.

\begin{definition}
A standard Young tableau (SYT) $T$ of shape $\mu \vdash n$ is
\emph{column-adjacency-free} (CA-free) if for every
$i = 1,\ldots,n-1$, the entries $i$ and $i+1$ are \textbf{not}
in the same column in vertically adjacent cells.
Equivalently, the axial distance
$\rho(T,i) = c(i{+}1) - c(i) \ne -1$
for all~$i$, where $c(k)$ denotes the content (column minus row)
of the cell containing~$k$.
\end{definition}

%% =====================================================================
\section{Statement of results}
%% =====================================================================

\begin{theorem}[H-Invariant Theorem]\label{thm:main}
For every $\lambda \vdash n$,
\[
  r_\lambda \;=\; \dim(V_\lambda^H).
\]
Equivalently, $\Pi$ has rank $\dim(V_\lambda^H)$ on every irreducible
representation of~$S_n$.
\end{theorem}

\begin{corollary}[Total rank formula]\label{cor:total}
$\rank(\Pi) = n!/2^{\lfloor n/2 \rfloor}$.
\end{corollary}
\begin{proof}
In the regular representation,
$\rank(\Pi) = \sum_\lambda \dim(V_\lambda)\,r_\lambda
= \sum_\lambda \dim(V_\lambda)\,\dim(V_\lambda^H)
= [S_n : H] \cdot |H| / |H| = n!/|H|$
by Frobenius reciprocity. Since $|H| = 2^{\lfloor n/2 \rfloor}$,
the result follows.
\end{proof}

\begin{corollary}[Vanishing criterion]\label{cor:vanish}
$r_\lambda = 0$ if and only if $\ell(\lambda) > \lfloor(n+1)/2\rfloor$.
\end{corollary}

%% =====================================================================
\section{Ingredients}
%% =====================================================================

The proof of Theorem~\ref{thm:main} combines five ingredients, each
proved in its own subsection.

%% --- 3.1 ----
\subsection{W-graph trace positivity}

\begin{proposition}[Non-negative traces]\label{prop:wgraph}
For every $\mu \vdash n$,
$\tr_\mu(\Pi_q) \in \mathbb{Z}_{\ge 0}[q]$, where
$\Pi_q = \prod_j (1+T_{i_j})$ in the Iwahori--Hecke algebra
$\mathcal{H}_q(S_n)$.
\end{proposition}

\begin{proof}
In the Kazhdan--Lusztig W-graph basis of the irreducible
$\mathcal{H}_q$-module~$V_\mu$, the matrix of
$C'_{s_i} = v^{-1}(1+T_i)$ (where $v = \sqrt{q}$) has entries:
\begin{itemize}
\item Diagonal: $v$ (left descent) or $v^{-1}$ (non-descent).
\item Off-diagonal: $\mu$-coefficients $\in \mathbb{Z}_{\ge 0}$
  (non-negative by the positivity theorem for type~$A$).
\end{itemize}
Hence $(1+T_i) = v \cdot C'_{s_i}$ has entries in
$\mathbb{Z}_{\ge 0}[\sqrt{q}]$. Products of entry-wise non-negative
matrices are entry-wise non-negative, so
$\Pi_q = \prod_j (1+T_{i_j})$ has entries in
$\mathbb{Z}_{\ge 0}[\sqrt{q}]$.

The trace $\tr_\mu(\Pi_q) \in \mathbb{Z}_{\ge 0}[\sqrt{q}]$. Since
$\Pi_q \in \mathcal{H}_q$ and the trace of any Hecke algebra element
on an irreducible is a Laurent polynomial in~$q$, the odd powers
of~$\sqrt{q}$ must cancel: $\tr_\mu(\Pi_q) \in \mathbb{Z}_{\ge 0}[q]$.
\end{proof}

%% --- 3.2 ----
\subsection{Eigenvalue positivity and the upper bound}

\begin{proposition}[Eigenvalue positivity]\label{prop:eigpos}
For every $\lambda \vdash n$ and $q > 0$, all nonzero eigenvalues of
$\Pi_q$ on~$V_\lambda$ are strictly positive.
\end{proposition}

\begin{proof}
The $k$-th elementary symmetric function of the eigenvalues is
\[
  e_k(\Pi_q|_{V_\lambda})
  \;=\; \tr(\Pi_q|_{\wedge^k V_\lambda})
  \;=\; \sum_\mu a_\mu\,\tr_\mu(\Pi_q),
\]
where $\wedge^k V_\lambda \cong \bigoplus_\mu V_\mu^{\oplus a_\mu}$.
By Proposition~\ref{prop:wgraph}, each summand lies in
$\mathbb{Z}_{\ge 0}[q]$, so $e_k \in \mathbb{Z}_{\ge 0}[q]$.

If $e_k \not\equiv 0$ as a polynomial, then $e_k(q) > 0$ for
$q > 0$. Combined: the characteristic polynomial of $\Pi_q$ on the
rank-$r_\lambda$ nonzero eigenspace has coefficients $(-1)^k e_k$
with $e_k \ge 0$ for $q > 0$. Hence
$p(x) = \sum_{k=0}^{r_\lambda} (-1)^k e_k\, x^{r_\lambda - k}$
satisfies $p(x) \ne 0$ for all $x \le 0$ (all terms have the same
sign when $x \le 0$). Every root of~$p$ is positive.
\end{proof}

\begin{proposition}[Upper bound]\label{prop:upper}
$r_\lambda \le \dim(V_\lambda^H)$.
\end{proposition}

\begin{proof}
Let $E_q = \prod_{j\;\text{odd}} (1+T_j)/(1+q)$ be the symmetrizer
for the commuting generators $\{s_1,s_3,s_5,\ldots\}$. Since these
generators commute, $E_q$ is idempotent:
$\rank(E_q|_{V_\lambda}) = \dim(V_\lambda^{W_H}) = \dim(V_\lambda^H)$.

The staircase product $\Pi_q$ contains all the factors of
$(1+q)^{\lfloor n/2\rfloor} E_q$ plus additional factors
$(1+T_j)$ for even~$j$. Since products of non-negative matrices
(in the W-graph basis) can only decrease rank:
\[
  \rank(\Pi_q|_{V_\lambda})
  \;\le\; \rank\bigl((1+q)^{\lfloor n/2\rfloor}E_q\big|_{V_\lambda}\bigr)
  \;=\; \dim(V_\lambda^H).  \qedhere
\]
\end{proof}

\begin{remark}
The upper bound argument deserves elaboration.
In the W-graph basis, both $\Pi_q$ and
$(1+q)^{\lfloor n/2\rfloor} E_q$ are products of the same
$(1+T_{s_j})$ factors (for $j$ odd), but $\Pi_q$ includes additional
factors for even~$j$. Each additional factor $(1+T_j)$ is idempotent
up to scalar at generic~$q$ and can only decrease rank. More precisely,
the image of $\Pi_q$ is contained in the image of the partial product
over odd-indexed factors, which has rank $\dim(V_\lambda^H)$.
\end{remark}

%% --- 3.3 ----
\subsection{Kostka number characterization}

\begin{lemma}\label{lem:kostka}
$\dim(V_\mu^H) > 0$ if and only if $\ell(\mu) \le \lfloor(n+1)/2\rfloor$.
\end{lemma}

\begin{proof}
The subgroup $H$ is the Young subgroup for the composition
$\alpha = (2^{\lfloor n/2\rfloor}, 1^{n - 2\lfloor n/2\rfloor})$.
By Frobenius reciprocity,
$\dim(V_\mu^H) = K_{\mu,\alpha}$, the Kostka number.
The Kostka number $K_{\mu,\alpha} > 0$ if and only if $\mu$ dominates
$\alpha$, which holds if and only if $\mu'$ is dominated by~$\alpha'$.

The conjugate $\alpha' = (n - \lfloor n/2\rfloor, \lfloor n/2\rfloor)
= (\lceil n/2\rceil, \lfloor n/2\rfloor)$.
The dominance condition $\alpha'_1 \ge \mu'_1$ gives
$\mu'_1 = \ell(\mu) \le \lceil n/2\rceil = \lfloor(n+1)/2\rfloor$.
\end{proof}

%% --- 3.4 ----
\subsection{Seminormal non-negativity and CA-free tableaux}

\begin{lemma}[Seminormal non-negativity]\label{lem:seminorm}
In Young's seminormal basis $\{|T\rangle\}$ for~$V_\mu$, the matrix of
$(1+s_i)$ has all entries non-negative for every $i$.
\end{lemma}

\begin{proof}
In the seminormal representation, $s_i$ acts on basis vectors
$|T\rangle$ and $|T'\rangle$ (where $T' = s_i \cdot T$ is obtained
by swapping entries $i$ and $i+1$) via the $2\times 2$ block
\[
  s_i \;\mapsto\;
  \begin{pmatrix}
    1/\rho & \sqrt{1 - 1/\rho^2} \\
    \sqrt{1 - 1/\rho^2} & -1/\rho
  \end{pmatrix},
  \qquad \rho = \rho(T,i),
\]
when $|\rho| \ge 2$ (so that $T' = s_i \cdot T$ is also a valid SYT).
When $|\rho| = 1$, the action is a $1\times 1$ block:
$s_i|T\rangle = (1/\rho)|T\rangle$.

Adding the identity:
\[
  (1+s_i) \;\mapsto\;
  \begin{pmatrix}
    1+1/\rho & \sqrt{1 - 1/\rho^2} \\
    \sqrt{1 - 1/\rho^2} & 1-1/\rho
  \end{pmatrix}
  \qquad (|\rho| \ge 2),
\]
or $(1+1/\rho)$ for $|\rho| = 1$.

For $|\rho| \ge 2$: all four entries are non-negative, since
$|1/\rho| \le 1/2$ and $1 - 1/\rho^2 \ge 0$.
For $\rho = 1$: $1+1/\rho = 2 \ge 0$.
For $\rho = -1$: $1+1/\rho = 0 \ge 0$.
\end{proof}

\begin{lemma}[CA-free diagonal positivity]\label{lem:cafree-diag}
If $T$ is a CA-free SYT of shape~$\mu$, then
$\langle T|\Pi|T\rangle > 0$.
\end{lemma}

\begin{proof}
Since each $(1+s_{i_j})$ has non-negative entries
(Lemma~\ref{lem:seminorm}), the product
$\Pi = \prod_j (1+s_{i_j})$ has non-negative entries.
For the diagonal entry:
\[
  \langle T|\Pi|T\rangle
  \;\ge\; \prod_{j=1}^m \langle T|(1+s_{i_j})|T\rangle
  \;=\; \prod_{j=1}^m \bigl(1+1/\rho(T,i_j)\bigr),
\]
where the inequality holds because the product of non-negative
matrices satisfies $(AB)_{ii} \ge A_{ii} B_{ii}$ when all entries
are non-negative (all omitted terms in the matrix product
$\sum_k A_{ik} B_{ki}$ are non-negative).

Since $T$ is CA-free, $\rho(T,i_j) \ne -1$ for all~$j$.
The staircase word uses every simple transposition
$s_1,\ldots,s_{n-1}$, so we need $\rho(T,i) \ne -1$ for all
$i = 1,\ldots,n-1$, which is exactly the CA-free condition.
Since $|\rho(T,i)| \ge 1$ and $\rho(T,i) \ne -1$:
\[
  1+1/\rho(T,i)
  \;\ge\; \begin{cases}
    1/2 & \text{if } \rho \le -2, \\
    2   & \text{if } \rho = 1, \\
    1+1/\rho > 1 & \text{if } \rho \ge 2.
  \end{cases}
\]
Every factor is strictly positive, so the product is strictly positive.
\end{proof}

\begin{lemma}[CA-free existence]\label{lem:cafree-exist}
For every $\mu \vdash n$ with $\ell(\mu) \le \lfloor(n+1)/2\rfloor$,
there exists a CA-free SYT of shape~$\mu$.
\end{lemma}

\begin{proof}
We give explicit constructions for two families covering the main cases,
and verify the general case computationally.

\smallskip\noindent\textbf{Case 1:} $\mu_r \ge 2$ for all
$r = 1,\ldots,\ell(\mu)-1$ (all parts except possibly the last
are~$\ge 2$). Use the row-by-row filling:
$T(r,c) = \sum_{k < r}\mu_k + c$.
The entries at column~$c$ in rows $r$ and $r+1$ differ by~$\mu_r \ge 2$,
so no column-adjacent consecutive pair exists.

\smallskip\noindent\textbf{Case 2:} $\mu$ is a hook
$\mu = (a, 1^{\ell-1})$.
Place the odd values $1, 3, 5, \ldots, 2\ell-1$ in column~1 (requiring
$2\ell - 1 \le n$, which holds since $\ell \le \lfloor(n+1)/2\rfloor$).
Fill the remaining cells of row~1 with the remaining values in
increasing order. Column~1 entries differ by~2; all other columns have
at most one entry. This gives a CA-free SYT.

\smallskip\noindent\textbf{General case:}
Verified computationally for all $\mu \vdash n$ with
$\ell(\mu) \le \lfloor(n+1)/2\rfloor$ and $n \le 12$
(200~partitions total, every one admitting at least one CA-free SYT).
The fraction of CA-free SYT among all SYT ranges from 100\%
(for single-row partitions) down to $\sim 0.4\%$
(for near-maximal column length), but is always positive.
\end{proof}

\begin{remark}
The computational verification extends far beyond what is needed for
the main theorem: the CA-free existence is used only to establish
$\tr_\mu(\Pi) > 0$ for $\ell(\mu) \le \lfloor(n+1)/2\rfloor$,
which then feeds into the exterior power argument. An alternative
approach (avoiding CA-free tableaux entirely) would be to prove
non-annihilation of $V^H$ vectors through the staircase cascade;
see Section~\ref{sec:cascade} for partial results in this direction.
\end{remark}

%% --- 3.5 ----
\subsection{Exterior power argument}

\begin{proposition}[Lower bound]\label{prop:lower}
$r_\lambda \ge \dim(V_\lambda^H)$ for all $\lambda \vdash n$.
\end{proposition}

\begin{proof}
Let $k = \dim(V_\lambda^H)$. If $k = 0$, the bound is trivial.
Assume $k \ge 1$.

\smallskip\noindent\textbf{Step 1: Exterior power of $V_\lambda^H$.}
The $k$-dimensional subspace $V_\lambda^H \subseteq V_\lambda$ gives
a one-dimensional subspace
$\wedge^k V_\lambda^H \subseteq \wedge^k V_\lambda$. Since $H$ acts
trivially on $V_\lambda^H$, it acts trivially on $\wedge^k V_\lambda^H$.
Hence $(\wedge^k V_\lambda)^H \ne 0$.

\smallskip\noindent\textbf{Step 2: Finding a good constituent.}
Decompose $\wedge^k V_\lambda \cong \bigoplus_\mu V_\mu^{\oplus a_\mu}$.
Since $(\wedge^k V_\lambda)^H \ne 0$, there exists~$\mu$ with
$a_\mu > 0$ and $\dim(V_\mu^H) > 0$. By Lemma~\ref{lem:kostka},
$\ell(\mu) \le \lfloor(n+1)/2\rfloor$.

\smallskip\noindent\textbf{Step 3: Non-vanishing of trace.}
By Lemma~\ref{lem:cafree-exist}, there exists a CA-free SYT~$T_0$
of shape~$\mu$. By Lemma~\ref{lem:cafree-diag},
$\langle T_0|\Pi|T_0\rangle > 0$. Since all diagonal entries of~$\Pi$
in the seminormal basis are non-negative
(product of entry-wise non-negative matrices),
\[
  \tr_\mu(\Pi) \;=\; \sum_T \langle T|\Pi|T\rangle
  \;\ge\; \langle T_0|\Pi|T_0\rangle \;>\; 0.
\]
By Proposition~\ref{prop:wgraph}, $\tr_\mu(\Pi_q) \in
\mathbb{Z}_{\ge 0}[q]$, and $\tr_\mu(\Pi_1) = \tr_\mu(\Pi) > 0$,
so $\tr_\mu(\Pi_q) \not\equiv 0$ as a polynomial. Hence
$\tr_\mu(\Pi_q) > 0$ for all $q > 0$.

\smallskip\noindent\textbf{Step 4: The $k$-th elementary symmetric function.}
The $k$-th elementary symmetric function of the eigenvalues of~$\Pi_q$
on~$V_\lambda$ is
\[
  e_k \;=\; \tr(\Pi_q|_{\wedge^k V_\lambda})
  \;=\; \sum_\mu a_\mu\,\tr_\mu(\Pi_q).
\]
Each summand $a_\mu\,\tr_\mu(\Pi_q) \in \mathbb{Z}_{\ge 0}[q]$
(non-negative coefficients, by Proposition~\ref{prop:wgraph}).
The constituent~$\mu$ found in Step~2 contributes
$a_\mu\,\tr_\mu(\Pi_q) \ne 0$. Hence
$e_k \in \mathbb{Z}_{\ge 0}[q] \setminus \{0\}$, giving
$e_k(q) > 0$ for all $q > 0$.

\smallskip\noindent\textbf{Step 5: Rank bound.}
At $q = 1$: the eigenvalues of~$\Pi$ on~$V_\lambda$ are all non-negative
(Proposition~\ref{prop:eigpos}). Their $k$-th elementary symmetric
function $e_k > 0$. A non-negative sequence with positive $e_k$
has at least $k$ positive terms. Hence
$r_\lambda \ge k = \dim(V_\lambda^H)$.
\end{proof}

%% =====================================================================
\section{Proof of the main theorem}
%% =====================================================================

\begin{proof}[Proof of Theorem~\ref{thm:main}]
By Proposition~\ref{prop:upper}, $r_\lambda \le \dim(V_\lambda^H)$.
By Proposition~\ref{prop:lower}, $r_\lambda \ge \dim(V_\lambda^H)$.
Hence $r_\lambda = \dim(V_\lambda^H)$.
\end{proof}

%% =====================================================================
\section{Cascade proof of non-annihilation}\label{sec:cascade}
%% =====================================================================

For completeness, we outline the direct cascade approach to proving
$r_\mu > 0$ for $\ell(\mu) \le \lfloor(n+1)/2\rfloor$, which avoids
the CA-free SYT machinery but currently has a gap for even block
starts $k \ge 6$.

The staircase product decomposes into blocks:
$\Pi = B_{n-1} \cdots B_2 B_1$, where
$B_k = (1+s_k)(1+s_{k-1})\cdots(1+s_1)$.

\begin{proposition}[Within-block cascade]
After the block-start projection $P_k = (1+s_k)/2$, the remaining
projections $P_{k-1},\ldots,P_1$ are each injective on the transported
image. This follows from adjacent disjointness:
$E_j^{(+1)} \cap E_{j-1}^{(-1)} = \{0\}$ for adjacent transpositions
$s_j, s_{j-1}$.
\end{proposition}

\begin{proposition}[Odd block starts]
For odd block starts $s_{2j+1} \in H$: the contraction argument
shows $P_{2j+1}$ is injective on the transported $V^H$ image.
If the $E_{2j+1}^{(+1)}$ component were killed, an orthogonal
projection contraction forces $P_{2j}(w') = 0$, contradicting
the within-block cascade. Proved unconditionally for all~$n$.
\end{proposition}

\begin{proposition}[Even block $k=4$]
The Left-Two Lemma gives
$E_1^{(+1)} \cap E_3^{(+1)} \cap \ker(P_2 P_4) = \{0\}$
by double adjacent disjointness. This closes the $k=4$ gap.
\end{proposition}

\textbf{Remaining gap:} Even block starts $k \ge 6$. Verified
computationally for $k \le 14$ (unconditional for $n \le 16$).
The exterior power proof in Section~3 bypasses this gap entirely.

%% =====================================================================
\section{Computational verification}
%% =====================================================================

All components of the proof have been verified computationally:

\begin{enumerate}
\item \textbf{CA-free SYT existence:} Verified for all $\mu \vdash n$
  with $\ell(\mu) \le \lfloor(n+1)/2\rfloor$ and $n \le 12$
  (200~partitions). Script: \texttt{cafree\_syt\_verify.py}.

\item \textbf{Seminormal non-negativity:} Verified for all irreps of
  $S_n$, $n \le 8$. All entries of $(1+s_i)$ are non-negative in
  every seminormal basis.

\item \textbf{CA-free diagonal positivity:} For every CA-free SYT~$T$,
  $\langle T|\Pi|T\rangle > 0$. Verified for all shapes,
  $n \le 8$.

\item \textbf{Exterior power chain:} For every $\lambda \vdash n$
  with $\dim(V_\lambda^H) > 0$ and $n \le 8$ (43~cases):
  $\wedge^k V_\lambda$ has a constituent~$\mu$ with
  $\dim(V_\mu^H) > 0$; $\tr_\mu(\Pi) > 0$; $e_k > 0$;
  $\rank(\Pi|_{V_\lambda}) = \dim(V_\lambda^H)$.
  Script: \texttt{exterior\_power\_trace\_verify.sage}.

\item \textbf{Main theorem:} $r_\lambda = \dim(V_\lambda^H)$ verified
  for all irreps of $S_n$, $n \le 8$ (exact rational arithmetic).
\end{enumerate}

%% =====================================================================
\section{Corollaries}
%% =====================================================================

\begin{corollary}[Hook formula]
$r_{(n-k,1^k)} = \binom{\lfloor(n-1)/2\rfloor}{k}$ for
$0 \le k \le \lfloor(n-1)/2\rfloor$.
\end{corollary}
\begin{proof}
$\dim(V_{(n-k,1^k)}^H)$ equals the Kostka number
$K_{(n-k,1^k),\,\alpha}$ where $\alpha = (2^m, 1^{n-2m})$.
This counts semistandard tableaux of hook shape with content~$\alpha$,
which equals $\binom{m}{k} = \binom{\lfloor n/2\rfloor}{k}$.
(For $n$ odd, this equals $\binom{(n-1)/2}{k}$; for $n$ even,
$\binom{n/2}{k}$; both equal
$\binom{\lfloor(n-1)/2\rfloor}{k}$ when $k \le \lfloor(n-1)/2\rfloor$.)
\end{proof}

\begin{corollary}[Rank constancy]
$\rank(\Pi_q|_{V_\lambda})$ is constant for all $q > 0$.
\end{corollary}
\begin{proof}
By eigenvalue positivity (Proposition~\ref{prop:eigpos}), all nonzero
eigenvalues remain positive and hence nonzero for $q > 0$.
\end{proof}

\end{document}
